\documentclass[11pt]{article}
\usepackage{amsmath,amssymb,amsthm}
\usepackage[margin=1in]{geometry}
\usepackage{mathtools}
\usepackage[hidelinks]{hyperref}

\newtheorem{theorem}{Theorem}
\newtheorem{lemma}[theorem]{Lemma}
\newtheorem{definition}[theorem]{Definition}
\newtheorem{corollary}[theorem]{Corollary}
\newtheorem{remark}[theorem]{Remark}

\newcommand{\pair}[2]{\langle #1, #2 \rangle}
\newcommand{\trueT}{\mathbf{t}}
\newcommand{\falseT}{\mathbf{f}}
\newcommand{\SchemaF}{\textbf{Schema~F}}
\newcommand{\nat}[1]{\overline{#1}}
\newcommand{\code}[1]{\ulcorner #1 \urcorner}
\newcommand{\reify}[1]{\widehat{#1}}
\newcommand{\thFun}{\mathrm{thFun}}
\newcommand{\thFunTFor}{\mathrm{thFunTFor}}
\newcommand{\cstf}{\mathrm{cstf}}
\newcommand{\TreeEq}{\mathrm{TreeEq}}
\newcommand{\IfLf}{\mathrm{IfLf}}

\title{G\"odel's Second Incompleteness Theorem\\
  in a Guard/Rose Equational System:\\
  Formalization in Agda}
\author{}
\date{}

\begin{document}
\maketitle

\section{Overview}

We formalize G\"odel's second incompleteness theorem in a purely
equational system combining Guard's syntax~\cite{guard1963}
(binary trees with typed unary and binary functors) and Rose's
proof style~\cite{rose1961} (no propositional connectives, only
equations and tree induction).  The formalization is in Agda with
\texttt{-{}-safe -{}-without-K -{}-exact-split}.  No postulates.

The system is the same as in the companion document
\texttt{guard-godelII-summary.tex}, which describes the formal
system and the G\"odel I proof.

\section{The formal system (summary)}

\subsection*{Trees and terms}
The ground data are finite binary trees:
$a, b ::= O \mid \pair{a}{b}$,
with $\trueT = O$ (``true'') and $\falseT = \pair{O}{O}$ (``false'').
Terms are built from trees, variables, and functor applications.
An \emph{equation} is a pair $t = u$ of terms.

\subsection*{Tree equality}
$\TreeEq$ is a binary functor (part of the system) that compares
two trees:
\[
  \TreeEq(a, b) = \begin{cases}
    \trueT & \text{if } a = b, \\
    \falseT & \text{if } a \neq b.
  \end{cases}
\]
Precisely: $\TreeEq(O, O) = \trueT$;
$\TreeEq(O, \pair{a}{b}) = \falseT$;
$\TreeEq(\pair{a}{b}, O) = \falseT$;
$\TreeEq(\pair{a_1}{a_2}, \pair{b_1}{b_2})
  = \IfLf(\TreeEq(a_1, b_1),\; \pair{\TreeEq(a_2, b_2)}{\falseT})$.
In particular $\TreeEq(t, t) = \trueT$ for all~$t$,
proved by \SchemaF{} (\texttt{treeEqSelf}).

\subsection*{Derivation rules}
Derivations are built from computation axioms (identity, projections,
composition, recursion, conditionals, $\TreeEq$), structural
rules (reflexivity, symmetry, transitivity, congruence),
variable substitution, a hypothesis rule, and \SchemaF{}
(uniqueness of tree recursion).

\subsection*{Consistency}
The system is \emph{consistent} (relative to hypothesis~$H$)
if there is no derivation of $\trueT = \falseT$ from~$H$.

\section{Proof encoding}

\subsection*{The code function}
Every term~$t$ is encoded as a tree $\code{t}$, and every
equation $t = u$ as $\code{t = u} = \pair{\code{t}}{\code{u}}$.
Every functor is similarly encoded.

\subsection*{The proof-encoding functor}
$\thFunTFor : \mathrm{Fun}_1$ is defined by tree recursion
($\mathrm{Rec}$) with a step function $\mathrm{thFunStep}$ that
dispatches on 26 proof-rule tags (0--25).  For each tag~$k$,
the corresponding case functor computes the equation code from
the proof data and recursive results.

\subsection*{Rose Theorem~14 (\texttt{thm14E})}
For every derivation $H \vdash t = u$, there exist trees
$p_a, p_b$ such that:
\begin{enumerate}
\item $\thFun(\pair{p_a}{p_b}) = \code{t = u}$ \quad (meta-level correctness),
\item $\vdash \thFunTFor(\pair{\reify{p_a}}{\reify{p_b}}) = \reify{\code{t = u}}$
  \quad (object-level correctness, polymorphic in hypothesis).
\end{enumerate}
The proof is by induction on the derivation, with 26~cases
(verified individually in the Nelson case files).

Component~(2) is \emph{polymorphic}: the derivation uses only
computation axioms and structural rules, never the hypothesis
rule.  This is crucial for G\"odel~II.

\section{G\"odel I (strengthened)}

\begin{theorem}[G\"odel~I --- \texttt{godelI} in \texttt{ProofEAny.agda}]
\label{thm:godel1}
For any hypothesis equation~$H$:
\[
  \mathit{Consistent}(H) \;\to\; H \vdash G \;\to\; \bot
\]
where $G$ is the G\"odel sentence (``no proof tree encodes
a derivation of this sentence'').
\end{theorem}

\noindent
This strengthens the original formulation by removing the
$\mathrm{ProofE}(H)$ assumption.  The key enabler is
$\mathrm{mkProofEAny}$, which constructs a proof encoding
for any equation using the trans-from-two-refls technique.

\section{G\"odel II}

\subsection*{The internal consistency statement}

The \emph{code} of an equation $t = u$ is the tree
$\code{t = u} = \pair{\code{t}}{\code{u}}$.  In particular,
the code of the contradiction $\trueT = \falseT$ is
\[
  \code{\bot} \;=\; \code{\trueT = \falseT}
  \;=\; \pair{\code{O}}{\code{\pair{O}{O}}}.
\]

\begin{definition}
The \emph{internal consistency statement} $\mathrm{Con}$ is the equation
\[
  \mathrm{Con} \;:\quad
  \TreeEq\bigl(\thFunTFor(x),\;\reify{\code{\bot}}\bigr) \;=\; \falseT
\]
where $x$ is a free variable.  Since $\TreeEq(a, b) = \falseT$
means $a \neq b$, the formula $\mathrm{Con}$ says:
\emph{for every proof tree~$x$, the equation encoded by~$x$
is not $\trueT = \falseT$} --- i.e., the system does not
prove $\trueT = \falseT$.
\end{definition}

\noindent
When the hypothesis is a trivially true equation such as
$\trueT = \trueT$, this corresponds exactly to Guard's
Theorem~14: ``$\mathrm{th}(y) \neq \text{``}0 = 1\text{''}$
is valid but unprovable in BRA''~\cite{guard1963}.

\subsection*{Statement}

\begin{theorem}[G\"odel~II, constructive form --- \texttt{conToBot} in \texttt{GodelIIFull.agda}]
\label{thm:godel2c}
For any hypothesis equation~$H$:
\[
  H \vdash \mathrm{Con} \;\to\;
  H \vdash \trueT = \falseT.
\]
\end{theorem}

\noindent
That is: if the system derives $\mathrm{Con}$,
then it derives $\trueT = \falseT$.  This is the direct constructive
content --- no consistency assumption, no Markov's principle.

\begin{corollary}[G\"odel~II, standard form --- \texttt{godelII} in \texttt{GodelIIFull.agda}]
\label{thm:godel2}
For any hypothesis equation~$H$:
\[
  \mathit{Consistent}(H) \;\to\;
  H \vdash \mathrm{Con} \;\to\; \bot.
\]
\end{corollary}

\noindent
The corollary follows by applying $\mathit{Consistent}(H)$
to the derivation of $\trueT = \falseT$ from
Theorem~\ref{thm:godel2c}.

\begin{remark}
The constructive form (Theorem~\ref{thm:godel2c}) is stronger
than the standard form (Corollary~\ref{thm:godel2}).
Classically they are equivalent by contraposition.
Constructively, the standard form alone gives only
$H \vdash \mathrm{Con} \to \lnot\lnot(H \vdash
\trueT = \falseT)$, and extracting the double negation
would require Markov's principle or
the Friedman--Dragalin $A$-translation.
However, our proof directly produces the derivation
$H \vdash \trueT = \falseT$ from $H \vdash \mathrm{Con}$
without any classical reasoning, since the construction of
the derivation of~$\trueT = \falseT$ does not use the consistency
assumption --- it is applied only at the final step to obtain~$\bot$.
\end{remark}

\subsection*{Proof}

The proof has three components.

\medskip\noindent
\textbf{Step 1: Pre-compute the encoding of the G\"odel~I argument.}

\noindent
Specialize the G\"odel~I proof to the hypothesis
$H_0 = G$ (the G\"odel sentence itself).
With this hypothesis, $\mathrm{ruleHyp}$ gives
$G \vdash G$, and the G\"odel~I argument
(Theorem~\ref{thm:godel1}) produces a derivation:
\[
  d_\bot : G \vdash \trueT = \falseT.
\]
Apply Theorem~14 (\texttt{thm14E}) to~$d_\bot$ to obtain
the proof encoding $\mathit{enc}_\bot$ and the correctness proof:
\[
  \mathit{corrBot} : \thFunTFor(\mathit{enc}_\bot) = \reify{\code{\bot}}.
\]
Crucially, $\mathit{corrBot}$ is \emph{polymorphic in the hypothesis}:
it is a derivation using only computation axioms and structural
rules, not the hypothesis rule.  Therefore it holds in any
hypothesis context, not just under~$G$.

The term $\mathit{enc}_\bot$ is a specific closed term computed
at the Agda meta-level.  It encodes the entire G\"odel~I
proof: the diagonal (\texttt{diagFTarget}), tree equality
self-comparison (\texttt{treeEqSelf}), variable instantiation,
and all 26~Nelson case verifications through \texttt{thm14E}.

\medskip\noindent
\textbf{Step 2: Instantiate the consistency hypothesis.}

\noindent
Given $H \vdash \mathrm{Con}$, apply variable
substitution (\texttt{ruleInst}) with $x := \mathit{enc}_\bot$:
\[
  H \vdash \TreeEq(\thFunTFor(\mathit{enc}_\bot),\;\reify{\code{\bot}}) = \falseT.
\]

\medskip\noindent
\textbf{Step 3: Derive the contradiction.}

\noindent
From $\mathit{corrBot}$ (polymorphic, instantiated at~$H$):
\[
  H \vdash \thFunTFor(\mathit{enc}_\bot) = \reify{\code{\bot}}.
\]
Replace $\thFunTFor(\mathit{enc}_\bot)$ with $\reify{\code{\bot}}$
by congruence:
\[
  H \vdash \TreeEq(\reify{\code{\bot}},\;\reify{\code{\bot}}) = \falseT.
\]
But $\TreeEq(t, t) = \trueT$ for all~$t$ (by \SchemaF{}):
\[
  H \vdash \TreeEq(\reify{\code{\bot}},\;\reify{\code{\bot}}) = \trueT.
\]
By transitivity: $H \vdash \trueT = \falseT$,
contradicting $\mathit{Consistent}(H)$.
\qed

\section{Discussion}

\subsection*{Comparison with Guard and Rose}

Guard~\cite{guard1963} proves G\"odel~II (his Theorem~14) by
explicitly constructing a functor~$h(x)$ within BRA such that
$\mathrm{th}(h(x)) = \text{``}0 = 1\text{''}$ whenever
$\mathrm{th}(x) = j$ (the G\"odel sentence code).  This uses
his Theorem~13 (internalization: $f(x) = y \supset
\mathrm{th}(D_f(x)) = \text{``}fx = y\text{''}$).

Rose~\cite{rose1961} proves G\"odel~II (his Theorem~18) by
applying his Theorem~16 (internalized Schema~F) to the function
$a(\mathrm{se}(\mathrm{nu}(N), N), \mathrm{th}(y))$, using the
consistency hypothesis $\mathrm{th}(x) \neq e(\mathrm{th}(z))$.

Our proof takes a different route: rather than constructing~$h$
as an object-level functor or applying internalized Schema~F,
we \emph{pre-compute} the encoding $\mathit{enc}_\bot$ at the
Agda meta-level from the concrete G\"odel~I derivation.
The polymorphicity of Theorem~14's correctness proof
($\mathit{corrBot}$ holds in any hypothesis context) is the
key property that makes this work.

\subsection*{The role of Theorem~14}

The entire non-trivial content of the proof resides in
Theorem~14 (\texttt{thm14E}).  The 26~Nelson case files verify
that $\thFunTFor$ correctly tracks each derivation rule.
The G\"odel~II proof applies Theorem~14 to the specific
derivation~$d_\bot$ (the G\"odel~I argument with hypothesis~$G$),
obtaining an encoding whose $\thFunTFor$ image is the
contradiction code.  The instantiation of $\mathrm{Con}$
at this encoding then produces the contradiction.

\subsection*{Formalization status}

All results compile with Agda~2.9.0 using
\texttt{-{}-safe -{}-without-K -{}-exact-split}.
No postulates.  The key files are:
\begin{itemize}
\item \texttt{GodelII.agda}: G\"odel~I (original, 168 lines) +
  \texttt{treeEqSelf} + \texttt{diagFTarget}
\item \texttt{ProofEAny.agda}: \texttt{mkProofEAny} +
  strengthened G\"odel~I (\texttt{godelI})
\item \texttt{Thm14E.agda}: Theorem~14, all 26 cases
\item \texttt{GodelIIFull.agda}: G\"odel~II
  (\texttt{godelII}, $\sim$90 lines)
\item \texttt{Nelson*.agda}: 26 individual case verifications
\end{itemize}

\begin{thebibliography}{9}
\bibitem{guard1963}
  J.~R.~Guard,
  \emph{Lecture Notes on Recursive Arithmetic},
  Princeton University, 1963.

\bibitem{rose1961}
  H.~E.~Rose,
  ``On the consistency and undecidability of recursive arithmetic,''
  \emph{Zeitschr.\ f.\ math.\ Logik und Grundlagen d.\ Math.},
  vol.~7, pp.~124--135, 1961.
\end{thebibliography}

\end{document}
